\begin{abcliste}
\abc Der Kasten ist isoliert, also ist $Q = 0$, und das Gas verrichtet keine Arbeit, $W = 0$. Nach dem ersten Hauptsatz ist deshalb
\begin{align}
 \Delta U &= Q + W = \result{0}
\end{align}
Die innere Energie bleibt gleich. Bei einem idealen Gas hängt $U$ nur von der Temperatur ab, die Temperatur bleibt also ebenfalls gleich.

\abc Die innere Energie bleibt auch beim realen Gas gleich, die Temperatur und damit die mittlere kinetische Energie der Teilchen nimmt aber ab. Die innere Energie eines realen Gases besteht deshalb nicht nur aus der kinetischen Energie der Teilchen, sondern enthält auch \emph{potentielle Energie der Anziehungskräfte zwischen den Teilchen}. Beim Ausdehnen wächst der mittlere Abstand der Teilchen; gegen die Anziehung muss dabei Arbeit verrichtet werden, die potentielle Energie nimmt auf Kosten der kinetischen Energie zu.

\abc Auch dann ist der Kasten als Ganzes isoliert, und es wird keine Arbeit nach aussen abgegeben: Die innere Energie des Gases bleibt gleich. Die Reibungsarbeit bremst die Strömung, die dabei entstehende Wärme bleibt aber im Gas. Nach dem Ausgleich ist der Endzustand derselbe wie ohne Düse (ideales Gas: gleiche Temperatur wie zu Beginn). Während des Überströmens ist das Gas in der Ausgangskammer kälter und in der Zielkammer wärmer, bis sich die Temperaturen ausgleichen.
\end{abcliste}
